\documentclass[11pt,a4paper]{article}
\usepackage[utf8]{inputenc}
\usepackage[T1]{fontenc}
\usepackage[brazil]{babel}
\usepackage[margin=2cm]{geometry}
\usepackage{booktabs}
\usepackage{array}
\usepackage{multirow}

\setlength{\parskip}{2pt}
\setlength{\parindent}{0pt}

\title{\vspace{-1.2cm}MAC0425 EP2: Busca Adversarial para o Pac-Man}
\author{\textbf{Lucas Aguiar Garcia} \qquad \textbf{NUSP 13672770}}
\date{\vspace{-0.4cm}}

\begin{document}
\maketitle
\vspace{-0.6cm}

\section{Implementação}
Foram implementados três agentes de busca adversarial com profundidade limitada em
\texttt{multiAgents.py}. O H-Minimax maximiza nos nós do Pac-Man (agente 0) e minimiza em
cada fantasma. O Alpha-Beta acrescenta poda com comparações \emph{estritas} (não poda em
empate, conforme exigido). O Expectimax trata os fantasmas como nós de chance, calculando a
média uniforme dos sucessores. Em todos, \texttt{self.depth} conta apenas as jogadas do Pac-Man (os
nós de fantasma não incrementam a profundidade); nós não-terminais no limite de profundidade são
avaliados como folhas pela função de avaliação.

\section{Metodologia}
Cada um dos três algoritmos foi testado nos quatro \emph{layouts} (\texttt{minimaxClassic},
\texttt{mediumClassic}, \texttt{openClassic}, \texttt{trappedClassic}), com dois tipos de fantasma
(aleatório padrão e adversarial \texttt{-g DirectionalGhost}), duas funções de avaliação
(padrão = \texttt{scoreEvaluationFunction} e better = \texttt{betterEvaluationFunction})
e profundidades 1 e 3. Cada caso reporta a média da pontuação final de 3 partidas com semente
fixa. DNF indica partida que não termina: com a avaliação padrão num mapa aberto o Pac-Man não
tem incentivo a comer, vagueia indefinidamente e nenhuma pontuação final é atingida.

\section{Resultados}

\begin{table}[h]
\centering
\footnotesize
\caption{Fantasmas aleatórios: pontuação média de 3 partidas.}
\begin{tabular}{l c rr rr rr}
\toprule
 & & \multicolumn{2}{c}{Minimax} & \multicolumn{2}{c}{Alpha-Beta} & \multicolumn{2}{c}{Expectimax}\\
\cmidrule(lr){3-4}\cmidrule(lr){5-6}\cmidrule(lr){7-8}
Layout & $d$ & pad & bet & pad & bet & pad & bet \\
\midrule
\multirow{2}{*}{minimaxClassic} & 1 & -158 & -160 & -158 & -160 & -158 & 177 \\
                                & 3 & -159 & -159 & -159 & -159 & 511 & 174 \\
\multirow{2}{*}{mediumClassic}  & 1 & -1107 & 1910 & -1107 & 1910 & -1107 & 1592 \\
                                & 3 & -1970 & 1341 & -1970 & 1341 & -1165 & 1915 \\
\multirow{2}{*}{openClassic}    & 1 & DNF & 1310 & DNF & 1310 & DNF & 1245 \\
                                & 3 & DNF & 1332 & DNF & 1332 & DNF & 1231 \\
\multirow{2}{*}{trappedClassic} & 1 & 531 & -158 & 531 & -158 & 531 & -158 \\
                                & 3 & -501 & -501 & -501 & -501 & -157 & -157 \\
\bottomrule
\end{tabular}
\end{table}

\vspace{-0.3cm}
\begin{table}[h]
\centering
\footnotesize
\caption{Fantasmas direcionais (\texttt{-g DirectionalGhost}): pontuação média de 3 partidas.}
\begin{tabular}{l c rr rr rr}
\toprule
 & & \multicolumn{2}{c}{Minimax} & \multicolumn{2}{c}{Alpha-Beta} & \multicolumn{2}{c}{Expectimax}\\
\cmidrule(lr){3-4}\cmidrule(lr){5-6}\cmidrule(lr){7-8}
Layout & $d$ & pad & bet & pad & bet & pad & bet \\
\midrule
\multirow{2}{*}{minimaxClassic} & 1 & -493 & -494 & -493 & -494 & -493 & -161 \\
                                & 3 & -497 & -160 & -497 & -160 & -163 & -164 \\
\multirow{2}{*}{mediumClassic}  & 1 & 307 & 326 & 307 & 326 & 307 & 326 \\
                                & 3 & 432 & 1263 & 432 & 1263 & 540 & 1766 \\
\multirow{2}{*}{openClassic}    & 1 & DNF & 1257 & DNF & 1257 & DNF & 1339 \\
                                & 3 & DNF & 1299 & DNF & 1299 & DNF & 1336 \\
\multirow{2}{*}{trappedClassic} & 1 & -158 & -158 & -158 & -158 & -158 & -158 \\
                                & 3 & -501 & -501 & -501 & -501 & -502 & -502 \\
\bottomrule
\end{tabular}
\end{table}

\section{Discussão}

Minimax e Alpha-Beta são idênticos em pontuação. Em todas as células as duas colunas coincidem:
a poda não altera a ação escolhida, apenas elimina expansões redundantes. A vantagem do Alpha-Beta é
exclusivamente computacional. Com a mesma semente, em \texttt{openClassic}+better com $d=3$ o Minimax
levou de 26 a 33\,s contra 7 a 8\,s do Alpha-Beta ($\approx 4\times$ mais rápido), e cerca de $2\times$ em
\texttt{mediumClassic}. Quanto maior a profundidade e o fator de ramificação, maior o ganho da poda.

A função de avaliação pesa mais que a profundidade. A avaliação padrão (apenas o \emph{score})
não guia o agente a comer: em \texttt{mediumClassic} produz pontuações fortemente negativas
($-1107$, $-1970$) e em \texttt{openClassic} a partida sequer termina (DNF), pois num mapa aberto o
agente apenas evita os fantasmas e oscila. A \texttt{better}, que premia a proximidade da comida e pune
a dos fantasmas, transforma esses casos em vitórias ($1910$, $1341$, ${\sim}1300$ em
\texttt{openClassic}). Aumentar a profundidade de 1 para 3 ajuda bem menos que trocar pad por better.

Expectimax supera Minimax contra fantasmas aleatórios. Como os fantasmas reais se movem ao
acaso, o modelo de chance do Expectimax é mais fiel que a hipótese de pior caso do Minimax, que fica
``paranoico''. Exemplos: \texttt{minimaxClassic}/aleatório/$d{=}3$, Expectimax faz 511 (3/3
vitórias) contra $-159$ do Minimax; \texttt{mediumClassic}/aleatório/better/$d{=}3$, 1915
contra $1341$. O caso mais ilustrativo é \texttt{trappedClassic}: com $d=3$ o Minimax vê a morte como
inevitável e ``desiste'' (corre para o fantasma a fim de minimizar a penalidade de tempo, $-501$, 0/3),
enquanto o Expectimax aposta que o fantasma aleatório pode se afastar e às vezes escapa ($-157$, 1/3).

Contra fantasmas direcionais a vantagem do Expectimax diminui. Quando os fantasmas de fato
perseguem, situação próxima da hipótese de pior caso do Minimax, os três algoritmos se aproximam
(Tabela 2). O Expectimax ainda lidera em mapas onde buscar comida importa mais que evitar o pior caso
(\texttt{mediumClassic}/$d{=}3$/better: $1766$ contra $1263$), mas a diferença é menor do que contra
fantasmas aleatórios.

Conclusão. Minimax e Alpha-Beta entregam o mesmo jogo, com o Alpha-Beta sendo a escolha
eficiente; o Expectimax é preferível quando o oponente não é ótimo (fantasmas aleatórios); e a qualidade
da função de avaliação é o fator decisivo de desempenho na busca com profundidade limitada.

\end{document}
